% EVAL FIXTURE — see b1.tex. Present so the census has both parts to report on;
% deliberately thin, because stage 6 must not be reached in this scenario.
\documentclass[11pt,a4paper]{article}
\usepackage[b2]{ercformatting}
\usepackage{ercreview}

\begin{document}
\maketitle
\startsectiona

\section{Project design}
Three work packages, one per objective, run in sequence.

\startsectionb

\subsection{WP1: Grain-scale flux}
\emph{Objective.} A calibrated grain-flux record at both field sites.
\emph{Rationale.} O1 cannot be tested without simultaneous records at the two scales,
and no existing dataset contains both.

\paragraph{Task 1.1: Instrument deployment.}
To obtain simultaneous records we deploy paired sensors at both sites. Sensors are
co-located within one metre so that the grain-flux and bar-migration series share a
reference frame. This yields the first cross-scale record, the input to WP2.

\subsection{WP2: Scale coupling}
\emph{Objective.} A quantified transfer function between the two scales.
\emph{Rationale.} O2 asks which forcing sets the coupling, which requires the coupling
itself to be measurable first.

\subsection{Risk assessment}
Storm loss of instruments is the dominant risk. Two sites are instrumented so that
loss at one still yields a single-site coupling estimate, which answers O1 in reduced
form. The risk is worth taking because no weaker deployment resolves the coupling.

\startappendix
\begin{fundingtable}
No funding & & & & & \\
\end{fundingtable}

\end{document}
